\section{Introduction}

Large language models (LLMs) have recently shown strong potential as reasoning engines in financial analysis, enabling multi-agent systems to coordinate over heterogeneous signals-price data, technical indicators, chart patterns, and news sentiment-within a unified pipeline. While frameworks such as QuantAgent \cite{xiong2025quantagent} establish a compelling blueprint for LLM-driven trading agents, they are designed for mature, English-language markets and do not generalize to the constraints of emerging economies.

Vietnam's equity market presents three challenges that require explicit adaptation. First, a predictive bar horizon and the assumptions of an executable cash-equity simulation are distinct; calendar gaps and instrument scope make them non-interchangeable. Second, the dominant financial news source, CafeF, publishes exclusively in Vietnamese, making English-centric sentiment pipelines inapplicable. Third, the shallow cross-sectional universe of HOSE, HNX, and UPCoM calls for alpha factor selection methods suited to single-stock Open, High, Low, Close, and Volume (OHLCV) data rather than the large-universe assumptions of standard factor libraries.

We present a five-agent LangGraph \cite{langgraph2024} pipeline that addresses each of these gaps. An Alpha Agent dynamically selects the top-five factors from an executable registry of exactly 85 candidates: five proprietary OHLCV formulas and 80 single-stock adaptations drawn from the WorldQuant 101 catalogue \cite{kakushadze2016formulaic}. Candidates are ranked by point-in-time Information Coefficient (IC), accuracy, and a Sharpe diagnostic. This accounting is important: the registry contains neither 87 candidates nor all 101 WorldQuant expressions. The factor values and ranking use OHLCV only. CafeF sentiment is supplied separately as textual context for qualitative LLM comparison. The primary configured scorer is a hosted ViSoBERT \cite{nguyen2023visobert} endpoint, with a deterministic lexicon fallback; the scorer and fallback status are recorded per article, and a run is attributed to the endpoint only when its included articles verify that path. Pattern and Trend agents apply vision-language analysis to candlestick and trendline charts, while a Decision Agent synthesises all reports into a binary LONG/SHORT direction forecast whose account action is resolved separately.

The primary objective of this paper is to evaluate whether the Alpha Agent improves directional forecasting relative to a no-alpha baseline. Indicator, Pattern, and Trend reports are generated once at each cutoff and cloned into both decision branches, providing paired common-support estimates without upstream variation inside a comparison. The nine-symbol benchmark reports directional, statistical, and compounded-account outcomes. A separate four-way shared-upstream experiment then evaluates Full, Alpha-only, Sentiment-only, and Baseline variants, allowing the factor and sentiment contributions to be distinguished rather than treating the Alpha bundle as indivisible.

The remainder of this paper is organised as follows. Section 2 reviews related work on LLM-based trading agents, Vietnamese NLP, and quantitative alpha research. Section 3 formalises the problem. Section 4 describes the system architecture. Section 5 presents the experimental setup. Section 6 reports backtest results. Section 7 discusses limitations, and Section 8 concludes.
